% Pista ana-25j-q1 · Anàlisi
% Procedència: PAU Catalunya 2025 · Convocatòria de juny · Sèrie 1 · Exercici 1
\documentclass[11pt,a4paper]{article}
\usepackage{pista}
\pistainfo{Catalunya 2025}{Convocatòria de juny}{Sèrie 1}

\begin{document}

\section*{\textcolor{blauPista}{Exercici 1}}
\addcontentsline{toc}{section}{Exercici 1}

\noindent Considereu la funció
\[
f(x) = \dfrac{x^2 - 2x}{x - 1}.
\]

\subsection*{\textit{a)} \normalfont Determineu els talls de la corba $y = f(x)$ amb els eixos de coordenades, i les equacions de les seves possibles asímptotes verticals, horitzontals i obliqües. \hspace*{\fill}\textnormal{[1~punt]}}

\pista{Segueix l'ordre de l'estudi d'una funció:
\begin{itemize}
\item \textit{Domini i asímptotes verticals.} Mira on s'anul·la el denominador i calcula-hi el límit de la funció $f$. Si és $+\infty$ o $-\infty$, hi ha una asímptota vertical.
\item \textit{Límits a l'infinit.} Compara el grau del numerador i el del denominador. Quan el del numerador és una unitat més gran, no hi ha asímptota horitzontal, sinó una asímptota obliqua $y = mx + n$, amb:
\[
m = \lim_{x\to\infty} \frac{f(x)}{x}, \qquad n = \lim_{x\to\infty} \bigl(f(x) - mx\bigr)
\]
Per calcular $n$, fes la resta $f(x) - mx$ amb denominador comú.
\item \textit{Talls amb els eixos de coordenades.}

Per trobar el punt de tall amb l'eix d'abscisses, resol l'equació $f(x) = 0$.

Per trobar el punt de tall amb l'eix d'ordenades, és més senzill, perquè el punt serà $Q\bigl(0, f(0)\bigr)$.
\end{itemize}}

\subsection*{\textit{b)} \normalfont Calculeu les equacions de les rectes tangents a la corba $y = f(x)$ en els punts $x = 0$ i $x = 2$. Aquestes dues rectes són paral·leles? Justifiqueu la resposta. \hspace*{\fill}\textnormal{[1~punt]}}

\pista{En primer lloc, calcula $f'(x)$ amb la regla de la derivada d'un quocient:
\[
\left(\frac{u}{v}\right)' = \frac{u' \cdot v - u \cdot v'}{v^{2}}
\]
A continuació, en el punt $x = 0$, el pendent de la recta tangent és $f'(0)$, i el punt és $\bigl(0, f(0)\bigr)$. Per tant:
\[
y - f(0) = f'(0)\,(x - 0)
\]
Fes el mateix amb el punt $x = 2$.

\medskip
Finalment, dues rectes són paral·leles si tenen el mateix pendent: compara $f'(0)$ i $f'(2)$, i escriu la conclusió.}

\subsection*{\textit{c)} \normalfont Hi ha algun punt on la recta tangent a $f(x)$ tingui pendent 1? En cas afirmatiu, trobeu-lo. \hspace*{\fill}\textnormal{[0,5~punts]}}

\pista{El pendent de la recta tangent en un punt és el valor de la derivada en aquell punt. Per tant, planteja l'equació $f'(x) = 1$, amb la derivada de l'apartat \textit{b)}.

\medskip
Per resoldre-la, multiplica els dos costats pel denominador i desenvolupa el quadrat. Si arribes a una igualtat impossible, com ara $3 = 5$, vol dir que l'equació no té solució: llavors, no hi ha cap punt així. En canvi, si trobes algun valor de la variable $x$, la teva resposta han de ser els punts corresponents, cadascun amb les dues coordenades.}

\vspace{0.6em}

\end{document}
